\section{Method}
\label{sec:method}

\subsection{Cards as frozen feature vectors}
\label{sec:features}

Every card is a fixed 775-dimensional vector computed from public
information alone, frozen before training.

\paragraph{Structured features (391-d).}
Parsed from the card's Scryfall record: mana value (scaled + bucketed
one-hot, 10), mana-cost pips including hybrid/Phyrexian/X structure (10),
colors (8), color identity (5), supertypes (3), card types (9), creature
subtypes (top-128 vocabulary + an unmatched count, 129), power / toughness
/ loyalty each as (scaled, missing, star) triples (9), rarity (5), keyword
abilities (166 slots + unmatched count, 167), layout including back-face
flags for double-faced cards (10), and text-shape features: length, line
count, and 24 regex flags over rules text for common functional categories
(removal, card draw, counterspells, tokens, ramp, \ldots) (26).
Double-faced and split cards take numeric blocks from the front face.
Subtype and keyword vocabularies are counted over \emph{training} sets
only; a registry gate refuses to fit vocabularies on MSH, the frozen test
set (\S\ref{sec:protocol}). That guarantee covers MSH specifically. It does
not mean the dev trio's cards were excluded from every vocabulary in the
release (\S\ref{sec:data}), so calling the dev trio ``zero-shot'' is a
claim about training picks, not about where the vocabulary came from. The
resulting featurizer manifest is content-hashed
(\texttt{\FeaturizerHash}) and released; held-out cards featurize through
the frozen manifest, with out-of-vocabulary subtypes and keywords absorbed
by the unmatched counters. Scalars are clipped to $[0,1]$; parse failures
set a missing indicator, never a silent zero.

\paragraph{Text embedding (384-d).}
The card's type line and rules text, embedded with a frozen
\texttt{bge-small-en-v1.5} sentence encoder~\citep{xiao2023bge}
($L_2$-normalised). The card's own name, and for legendary cards its
short name, is masked to a placeholder token before embedding, so the
encoder cannot key on names it may know from pretraining. This is a
deliberate handicap, and it matters for licensed-IP sets
(\S\ref{sec:analysis}). Parenthetical reminder text is kept, since for a
brand-new mechanic it is the only definition the model will ever see. Back faces are appended for
double-faced cards.

\subsection{Architecture}
\label{sec:arch}

DraftFM is a two-tower scorer with \FdevParams M parameters (Figure~\ref{fig:arch}). Its two design
invariants operationalize the zero-shot contract (\S\ref{sec:data}):

\begin{enumerate}
  \item \textbf{No set identity.} There is no set embedding and no
  per-card parameter anywhere. The only way the model can know what set it
  is drafting is a \emph{set-context tower} that reads the set's full card
  list, plus the drafter's accumulating pool. ``Blue is the aggressive
  color in this set'' must be inferred from the card list, not memorized
  as a set identifier. Whether the tower itself turns out to be necessary
  for that inference, or whether the invariant holds just as well without
  it, is addressed later in this section.
  \item \textbf{Skill conditions the scorer, not the cards.} Card
  semantics are skill-invariant; the pick policy is not. Skill-bucket
  embeddings enter only the context pathway. Serving conditions on a fixed
  high-skill bucket; training sees every pick tagged with its drafter's
  bucket.
\end{enumerate}

\begin{figure}[t]
\centering
% Architecture diagram for the two-tower model (sections/method.tex, S4.2).
% Pure TikZ, no external assets -- compiles standalone as part of main.tex.
% \input from within a figure environment; see sections/method.tex.
\definecolor{archbox}{HTML}{F2F3F5}
\definecolor{archline}{HTML}{5B6472}
\definecolor{archblue}{HTML}{2563B0}
\definecolor{archmuted}{HTML}{6B7280}

\resizebox{\linewidth}{!}{%
\begin{tikzpicture}[
  font=\small,
  node distance=6mm and 8mm,
  archbox/.style={draw=archline, thick, rounded corners=1.5pt, fill=archbox,
              minimum height=8mm, align=center, inner sep=1.8mm},
  archtower/.style={draw=archblue, thick, rounded corners=1.5pt, fill=archblue!8,
              minimum height=9mm, align=center, inner sep=1.8mm},
  archarr/.style={-{Latex[length=2mm]}, thick, draw=archline},
]
  \node[archbox] (feat) {card feature\\vector\\\scriptsize 775-d};
  \node[archbox, right=of feat] (enc) {shared card\\encoder\\\scriptsize LayerNorm$\to$775--512--256, GELU};
  \node[archbox, right=of enc] (emb) {card\\embeddings\\\scriptsize 256-d};

  \node[archtower, above right=9mm and 9mm of emb] (ctx) {set-context tower\\\scriptsize cross-attn + rarity emb};
  \node[archtower, below right=9mm and 9mm of emb] (pool) {pool tower\\\scriptsize cross-attn, 4 queries};

  \node[archbox, above=7mm of ctx] (ctxin) {full set card list\\\scriptsize $\sim$400 rows};
  \node[archbox, below=7mm of pool] (poolin) {drafter's pool\\\scriptsize distinct cards + count emb};

  \node[archbox, right=15mm of ctx, yshift=-9mm] (cmlp) {context MLP\\\scriptsize +skill, format,\\pick-position feats};
  \node[archbox, right=10mm of cmlp] (score) {scorer MLP\\\scriptsize $[c;\,p;\,c{\odot}p;\,\mathrm{ctx}]$};
  \node[archbox, right=9mm of score] (softmax) {softmax\\\scriptsize over pack};

  \draw[archarr] (feat) -- (enc);
  \draw[archarr] (enc) -- (emb);
  \draw[archarr] (emb.north) .. controls +(0.5,0.3) and +(-0.7,0.1) .. (ctx.west);
  \draw[archarr] (emb.south) .. controls +(0.5,-0.3) and +(-0.7,-0.1) .. (pool.west);
  \draw[archarr] (ctxin) -- (ctx);
  \draw[archarr] (poolin) -- (pool);
  \draw[archarr] (ctx.east) -- (cmlp.north -| ctx.east) -- (cmlp.north);
  \draw[archarr] (cmlp) -- (score);
  \draw[archarr] (pool.east) .. controls +(1.0,0) and +(-1.0,-0.6) .. (score.south);
  \draw[archarr] (emb.south) .. controls +(0.4,-2.4) and +(-2.8,-1.7) .. (score.south)
    node[pos=0.24, above, font=\scriptsize, text=archmuted] {candidate $c$};
  \draw[archarr] (score) -- (softmax);
\end{tikzpicture}%
}

\caption{The two-tower scorer (\S\ref{sec:arch}). A shared card encoder
embeds every card once per (set, format) step; a pool tower summarises
the drafter's held cards, a set-context tower summarises the full card
list, and a per-candidate scorer combines the candidate embedding
$c$, the pool summary $p$, their elementwise product, and the context
vector into a softmax over the pack.}
\label{fig:arch}
\end{figure}

A shared \emph{card encoder} (LayerNorm $\to$ 775--512--256 MLP, GELU)
maps feature vectors to $d{=}256$ embeddings. Because training batches are
homogeneous in (set, format), the encoder runs once per step over the
set's full card list ($\sim$400 rows) and everything downstream gathers
rows from that table. The \emph{pool tower} represents the drafter's pool
as the set of distinct cards held, each embedding translated by a learned
count embedding (counts capped at 8), pooled by cross-attention from four
learned queries (4 heads); an empty pool attends over a learned null token.
The \emph{set-context tower} applies the same query-pooling to the entire
card list (with a rarity embedding added) to produce the set summary. A
context MLP combines the set summary with skill and format embeddings,
seven pick-position features (pack/pick indices, pool size), and four
set-shape scalars (card-list size, picks-per-pack indicator) into a 64-d
context. Each candidate $c$ in the pack is scored by an MLP over
$[\,c;\; p;\; c \odot p;\; \mathrm{ctx}\,]$, where $p$ is the pool summary
and $c \odot p$ the elementwise interaction; a softmax over the pack's
real slots gives the pick distribution.

We test the set-context tower's contribution directly with an experiment
that removes it and retrains; we call this variant A-noctx
(Table~\ref{tab:ablations}, \S\ref{sec:analysis}). It scores higher on all
three dev sets. That tower-free configuration, not the one just described,
is what F-full and the frozen MSH evaluation actually run. The first invariant still holds
either way, since neither configuration has a set-identity parameter.
Dropping the tower simply means the shipped model's only route to
set-relative signal is the pool tower, each candidate's own features, and
the four set-shape scalars above; it never reads the rest of the card
list.

\subsection{Training}
\label{sec:training}

One optimization recipe governs every run in the paper. It was fixed before the
scaling and variant studies, since the protocol forbids per-rung tuning.
Batches of 8{,}192 picks are drawn from one (set, format) shard at a
time. Shards are sampled proportionally to $n_{\text{picks}}^{0.5}$, an
interpolation between pick-uniform sampling, which lets the largest sets
dominate, and set-uniform sampling, which oversamples tiny sets. The
optimizer is AdamW ($\beta_1{=}0.9$, $\beta_2{=}0.98$, weight decay 0.01)
with peak learning rate $10^{-3}$, a 2{,}000-step linear warmup, and
cosine decay; the loss is cross-entropy with label smoothing 0.05,
renormalised over real pack slots. The presentation budget
is four epochs. Within-training validation (never the dev sets) is
evaluated every 2{,}000 steps and training stops after three
non-improving evaluations. Seed 17 throughout. Seed variance is measured
at two scaling rungs (\S\ref{sec:protocol}).

Training runs on a single Apple M3 Ultra workstation (96\,GB unified
memory) in PyTorch 2.12 eager mode on the MPS backend, at
$\sim$\FdevExamplesPerSec{} picks/s. The full \FdevPicks M-pick run takes
about \FdevWallHours{} hours, stopping at step 34{,}000. Training-backend
correctness is gated rather than assumed; specifics are in the companion
reference. Inference exports to ONNX (opset 18) and runs on a consumer
x86-64 machine, CPU-only, well under a millisecond per pick.
